% !TEX program = xelatex
% MetaBank 操作手册 - 使用 XeLaTeX 编译
\documentclass[12pt,a4paper]{article}
\usepackage{xeCJK}
\usepackage{amsmath,amssymb,amsfonts}
\usepackage{geometry}
\usepackage{graphicx}
\usepackage{hyperref}
\usepackage{booktabs}
\usepackage{listings}
\usepackage{xcolor}
\usepackage{titlesec}
\usepackage{fancyhdr}

\geometry{left=2.5cm,right=2.5cm,top=2.5cm,bottom=2.5cm}
\setCJKmainfont{SimSun}[AutoFakeBold=true]
\hypersetup{colorlinks=true,linkcolor=blue,urlcolor=blue}

\title{\textbf{MetaBank 元宇宙金融养老社区}\\[0.5em]\large 操作手册与技术原理}
\author{FIT「国脉杯」软件设计大赛 · 队伍：冲冲冲}
\date{\today}

\begin{document}
\maketitle

\begin{center}
\fbox{\parbox{0.9\textwidth}{
\textbf{文档说明：} 本文档使用 \textbf{LaTeX}（XeLaTeX 引擎）编写，包含数学公式、集合符号（如 $\mathcal{S}$, $\mathcal{F}$）、希腊字母（如 $\varepsilon$, $\sigma$）、下标上标等特殊符号。阅读 PDF 可获得最佳排版效果。
}}
\end{center}
\vspace{1em}

\tableofcontents
\newpage

%==============================================================================
\section{需求背景}
%==============================================================================

\subsection{老龄化社会的金融服务挑战}

根据国家统计局公开数据，截至 2024 年末，我国 60 岁及以上人口已达 2.97 亿，占总人口的 21.1\%，正加速迈入中度老龄化社会。与此同时，金融服务的数字化进程与老年群体的数字接纳能力之间形成了显著的结构性矛盾，具体表现在三个方面：

\begin{itemize}
\item \textbf{数字鸿沟。}主流金融应用以年轻用户为默认设计对象，字号偏小、交互层级深、术语密集。老年人上手困难，遇到风控弹窗、短信验证、人脸识别等环节时极易中断操作。
\item \textbf{金融诈骗高发。}老年群体因信息不对称与信任容错高，长期位列电信诈骗受害者榜首。传统金融平台缺乏"家庭监护 + 大额二次确认"机制，一次操作即可造成不可逆的资金损失。
\item \textbf{理财需求被忽视。}现有投顾产品以高风险高收益标的为主推，缺少针对老年人"稳健为先、通俗易懂"的专属路径。激进推荐反而会加剧风险敞口。
\end{itemize}

\subsection{项目定位}

MetaBank 以"金融养老社区"为切入点，将\textbf{元宇宙交互范式}（后续通过硬件穿戴设备在真实空间中实现）与\textbf{AI 智能顾问}结合，重塑老年金融服务体验。本系统的设计目标概括为一个三元组：
\begin{equation}
\mathcal{G} = (\text{可达性 Accessibility},\; \text{可信性 Trustworthiness},\; \text{可理解性 Comprehensibility})
\end{equation}
即任何年龄段用户均可便捷接入（可达），大额与风险操作有监护与拦截（可信），所有建议与状态都能以通俗语言被用户所理解（可理解）。

%==============================================================================
\section{目标用户画像与应用场景}
%==============================================================================

\subsection{典型用户画像}

本系统主要服务 60 岁以上银发群体及其家庭成员。通过用户访谈与画像聚类，我们提炼出三类典型角色：

\begin{table}[h]
\centering
\begin{tabular}{@{}llll@{}}
\toprule
\textbf{画像} & \textbf{年龄} & \textbf{特征} & \textbf{核心诉求} \\
\midrule
退休教师张大爷 & 70 & 有一定理财经验，视力下降 & 稳健保值、防诈骗 \\
独居李阿姨   & 65 & 不熟悉智能设备，子女在外 & 简单操作、子女可见 \\
银发创客王爷爷 & 68 & 对新科技有探索欲 & 元宇宙体验、社交互动 \\
\bottomrule
\end{tabular}
\caption{MetaBank 目标用户画像}
\end{table}

\subsection{典型应用场景}

\subsubsection{场景一：稳健理财}

张大爷登录 MetaBank 后，AI 顾问读取其风险偏好（保守型），在响应其"能不能多赚点"的询问时，主动规避高杠杆的国脉币期货 (GMFT)，转而推荐国脉币 (GMC) 稳定币定投与 AI 金融 ETF (AIFIN) 组合。系统在界面上以大字号、对比色强调"本金安全"与"预期年化"两个核心指标。

\subsubsection{场景二：家庭监护}

李阿姨在元宇宙商城购买"智能手环"等超出日常消费阈值的商品时，后端监护中间件根据 \texttt{daily\_limit} 配置触发二次确认流程：交易挂起并通过消息通道通知其已绑定的子女账户，由子女在移动端一键批准后方可落账。该流程有效缓解针对老年人的诱导性消费与诈骗。

\subsubsection{场景三：社区互动}

王爷爷通过穿戴式设备进入虚拟候鸟社区，在"社区公告牌"前浏览健康讲座通知，在"AI 顾问亭"前咨询近期市场，在"元宇宙商城"摊位挑选商品。整个交互以身体的走近与凝视触发，降低了按键与输入的认知负担。

\subsection{与现有系统的差异化}

\begin{itemize}
\item 面向\textbf{老年群体}进行专属设计，而非通用金融平台的老年模式。
\item \textbf{元宇宙交互}通过硬件穿戴设备在真实空间实现，避开纯 2D UI 的认知鸿沟。
\item \textbf{价格模型}采用几何布朗运动 + 跳跃扩散 + 宏观因子，对老年用户展示更符合真实市场逻辑的行情（详见第 4 节）。
\item \textbf{AI 顾问}以稳健保守为先验，主动拒绝高杠杆推荐。
\end{itemize}

%==============================================================================
\section{系统概述}
%==============================================================================

MetaBank 是基于元宇宙技术构建的金融养老社区平台，旨在解决老龄化群体在数字金融服务中的交互痛点。系统采用前后端分离架构，后端使用 Python FastAPI，前端使用 React + Vite，数据采用 JSON 文件存储。

\subsection{系统架构}

系统整体架构可抽象为：
\begin{equation}
\mathcal{S} = (\mathcal{F}, \mathcal{B}, \mathcal{D}, \mathcal{A})
\end{equation}
其中 $\mathcal{F}$ 为前端模块，$\mathcal{B}$ 为后端 API，$\mathcal{D}$ 为数据层，$\mathcal{A}$ 为 AI 智能助手模块。

各模块间的数据流满足：
\begin{equation}
\mathcal{F} \xrightarrow{\text{HTTP}} \mathcal{B} \xrightarrow{\text{读写}} \mathcal{D}, \quad \mathcal{A} \subset \mathcal{B}
\end{equation}

%==============================================================================
\section{认证与安全原理}
%==============================================================================

\subsection{JWT 认证机制}

系统采用 JWT (JSON Web Token) 进行无状态认证。令牌生成与验证遵循以下流程。

\subsubsection{令牌生成}

设用户标识为 $u$，过期时间为 $t_{\exp}$，密钥为 $K$，则载荷 (payload) 为：
\begin{equation}
P = \{ \mathtt{sub}: u,\; \mathtt{uid}: \mathrm{id}(u),\; \mathtt{exp}: t_{\exp},\; \mathtt{is\_admin}: \phi(u) \}
\end{equation}
其中 $\mathrm{id}(u)$ 为用户 ID，$\phi(u) \in \{0,1\}$ 表示是否为管理员。

JWT 由三部分组成，以点号连接：
\begin{equation}
\mathrm{JWT} = \mathrm{Base64URL}(\mathrm{Header}) \,\|\, \mathrm{Base64URL}(\mathrm{Payload}) \,\|\, \mathrm{Base64URL}(\mathrm{Signature})
\end{equation}

签名采用 HMAC-SHA256：
\begin{equation}
\mathrm{Signature} = \mathrm{HMAC}\text{-}\mathrm{SHA256}\bigl( \mathrm{Header} \| \mathrm{Payload},\; K \bigr)
\end{equation}

\subsubsection{令牌验证}

验证时需满足：
\begin{equation}
\mathrm{Verify}(\mathrm{JWT}, K) \land (t_{\mathrm{now}} < t_{\exp}) \implies \text{认证通过}
\end{equation}

\subsection{密码哈希}

密码采用 bcrypt 算法存储，不可逆。设明文密码为 $p$，盐值为 $s$，成本因子为 $c$，则：
\begin{equation}
h = \mathrm{bcrypt}(p, s, c) = \mathrm{EksBlowfish}\bigl( \mathrm{ExpandKey}(p, s, 2^c) \bigr)
\end{equation}

验证时：
\begin{equation}
\mathrm{verify}(p, h) \iff \mathrm{bcrypt}(p, h.s, h.c) = h
\end{equation}

\subsection{钱包地址生成}

用户钱包地址格式为 $\mathtt{0xGM} + H_{36}$，其中 $H_{36}$ 为 36 位十六进制串。生成算法：
\begin{equation}
r \sim \mathrm{Uniform}(\{0,1\}^{160}), \quad H = \mathrm{SHA256}(r)[:36]
\end{equation}
即先生成 20 字节随机数 $r$，再对其做 SHA-256 哈希并取前 36 个十六进制字符。

%==============================================================================
\section{交易所与定价原理}
%==============================================================================

为体现真实金融市场的运行逻辑，MetaBank 在 v2 中放弃了简单的随机扰动模型，转而采用\textbf{宏观因子驱动的几何布朗运动}（GBM）叠加 \textbf{Merton 跳跃扩散}的混合定价框架。本节给出完整的数学表述与系统实现。

\subsection{标的与基准价}

设标的集合为 $\mathcal{T} = \{\tau_1, \tau_2, \ldots, \tau_n\}$，每个标的 $\tau_i$ 有基准价 $P_i^{(0)}$：
\begin{equation}
P_{\mathtt{GM}}^{(0)} = 12.5,\;\; P_{\mathtt{GMAI}}^{(0)} = 25.0,\;\; P_{\mathtt{GMC}}^{(0)} = 1.0,\;\; P_{\mathtt{GMFT}}^{(0)} = 1.05,\;\; P_{\mathtt{METAV}}^{(0)} = 100.0,\;\; P_{\mathtt{AIFIN}}^{(0)} = 50.0
\end{equation}

\subsection{宏观因子建模}
\label{sec:macro}

引入四个宏观因子刻画市场环境，记为向量 $\mathbf{F}_t = (\mathrm{CPI}_t,\, \mathrm{AI}_t,\, \mathrm{Gold}_t,\, \mathrm{Sent}_t)^{\top}$，其演化分别为：

\begin{itemize}
\item \textbf{通胀指数 CPI}（随机游走，钳制于 $[0.5\%, 5\%]$）：
\begin{equation}
\mathrm{CPI}_{t+1} = \mathrm{clip}\bigl(\mathrm{CPI}_t + \eta^{c}_t,\; 0.5,\; 5.0\bigr),\quad \eta^{c}_t \sim \mathcal{N}(0, 0.05^2)
\end{equation}

\item \textbf{AI 热度}（Ornstein--Uhlenbeck 均值回复过程，$\theta=0.15$，$\mu_{\mathrm{AI}}=60$，钳制于 $[0,100]$）：
\begin{equation}
\mathrm{AI}_{t+1} = \mathrm{AI}_t + \theta\,(\mu_{\mathrm{AI}} - \mathrm{AI}_t) + \eta^{a}_t,\quad \eta^{a}_t \sim \mathcal{N}(0, 8^2)
\end{equation}

\item \textbf{黄金溢价}（与 CPI 正相关，钳制于 $[-2\%, 3\%]$）：
\begin{equation}
\mathrm{Gold}_{t+1} = 0.5\,(\mathrm{CPI}_{t+1} - 2.0) + \eta^{g}_t,\quad \eta^{g}_t \sim \mathcal{N}(0, 0.4^2)
\end{equation}

\item \textbf{市场情绪}（一阶自回归 AR(1)，钳制于 $[-1, +1]$）：
\begin{equation}
\mathrm{Sent}_{t+1} = 0.7\,\mathrm{Sent}_t + \eta^{s}_t,\quad \eta^{s}_t \sim \mathcal{N}(0, 0.25^2)
\end{equation}
\end{itemize}

为节省计算开销，宏观因子在内存中按 60 秒 TTL 缓存，并通过 WebSocket 通道每 5 秒推送至前端宏观面板。

\subsection{资产画像与差异化参数}
\label{sec:profile}

每个标的 $\tau$ 配备独立的画像参数 $\Theta_\tau = (\mu^{\tau}_{\mathrm{base}},\, \sigma^{\tau},\, \lambda^{\tau},\, \mu^{\tau}_J,\, \sigma^{\tau}_J,\, \boldsymbol{\beta}^{\tau})$，其中 $\boldsymbol{\beta}^{\tau} = (\beta^{\tau}_{\mathrm{cpi}},\, \beta^{\tau}_{\mathrm{ai}},\, \beta^{\tau}_{\mathrm{gold}},\, \beta^{\tau}_{\mathrm{sent}})$ 为四宏观因子敏感度。具体取值见下表：

\begin{table}[h]
\centering
\setlength{\tabcolsep}{4pt}
\begin{tabular}{@{}lcccccccc@{}}
\toprule
\textbf{标的} & $\boldsymbol{\mu_{\text{base}}}$ & $\boldsymbol{\sigma}$ & $\boldsymbol{\lambda}$ & $\boldsymbol{\sigma_J}$ & $\boldsymbol{\beta_{\text{cpi}}}$ & $\boldsymbol{\beta_{\text{ai}}}$ & $\boldsymbol{\beta_{\text{gold}}}$ & $\boldsymbol{\beta_{\text{sent}}}$ \\
\midrule
GM (国脉科技)     & 0.08 & 0.25 & 6  & 0.04 & 0.3  & 0.5 & 0.0 & 0.4 \\
GMAI (AI 概念)    & 0.15 & 0.40 & 12 & 0.07 & 0.2  & 1.2 & 0.0 & 0.6 \\
GMC (国脉币)      & 0.02 & 0.05 & 1  & 0.01 & -0.1 & 0.0 & 0.0 & 0.1 \\
GMFT (期货)       & 0.10 & 0.45 & 18 & 0.06 & 0.0  & 0.4 & 0.0 & 0.8 \\
METAV (元宇宙指数)& 0.12 & 0.30 & 8  & 0.05 & 0.2  & 0.8 & 0.0 & 0.5 \\
AIFIN (AI 金融 ETF)& 0.10 & 0.22 & 5  & 0.035& 0.1  & 0.7 & 0.0 & 0.4 \\
\bottomrule
\end{tabular}
\caption{六个标的的差异化模型参数（年化口径）}
\end{table}

参数取值的业务直觉：稳定币 GMC 给予最低 $\sigma$ 与 $\lambda$，期货 GMFT 给予最高 $\lambda$ 与 $\beta_{\mathrm{sent}}$，AI 概念 GMAI 给予最高 $\beta_{\mathrm{ai}}$。

\subsection{有效漂移率}

价格漂移由"基础漂移 + 宏观因子线性组合"得到：
\begin{align}
\mu^{\tau}_{\mathrm{eff}}(t) ={}& \mu^{\tau}_{\mathrm{base}} + \beta^{\tau}_{\mathrm{cpi}} \cdot \frac{\mathrm{CPI}_t - 2}{100} + \beta^{\tau}_{\mathrm{ai}} \cdot \frac{\mathrm{AI}_t - 60}{40} \cdot 0.1 \notag \\
& + \beta^{\tau}_{\mathrm{gold}} \cdot \frac{\mathrm{Gold}_t}{100} + \beta^{\tau}_{\mathrm{sent}} \cdot \mathrm{Sent}_t \cdot 0.05
\end{align}
其中 AI 热度与情绪系数的归一化常数（$0.1$ 与 $0.05$）使各因子对漂移率的边际贡献处于同一数量级。

\subsection{几何布朗运动 + Merton 跳跃扩散}

价格的连续时间动力学满足以下随机微分方程：
\begin{equation}
\boxed{\;dS^{\tau}_t = \mu^{\tau}_{\mathrm{eff}}(t)\, S^{\tau}_t\, dt + \sigma^{\tau} S^{\tau}_t\, dW_t + \bigl(J^{\tau}_t - 1\bigr) S^{\tau}_t\, dN_t\;}
\end{equation}
其中：
\begin{itemize}
\item $W_t$ 为标准 Brown 运动，提供连续扩散项；
\item $N_t \sim \mathrm{Poisson}(\lambda^{\tau} \cdot t)$ 为跳跃计数过程；
\item $\ln J^{\tau}_t \sim \mathcal{N}(\mu^{\tau}_J,\, (\sigma^{\tau}_J)^2)$ 为对数正态跳跃幅度。
\end{itemize}

跳跃项捕捉了真实市场中由突发事件（监管、产业政策、AI 热点、黑天鹅）触发的价格不连续。期货标的 GMFT 配置较高的 $\lambda$ 与 $\sigma_J$，使其呈现典型的"高频厚尾"特征；稳定币 GMC 则几乎平稳运行。

\subsection{离散日级递推}

设步长 $\Delta t = 1/252$（一个交易日），运用 Itô 引理对 $\ln S_t$ 离散化，可得：
\begin{equation}
S^{\tau}_{t+\Delta t} = S^{\tau}_t \cdot \exp\!\left[\bigl(\mu^{\tau}_{\mathrm{eff}} - \tfrac{1}{2}(\sigma^{\tau})^2\bigr)\Delta t + \sigma^{\tau}\sqrt{\Delta t}\,Z_t + \mathbb{1}_{U < \lambda^{\tau} \Delta t} \cdot \xi_t\right]
\end{equation}
其中 $Z_t \sim \mathcal{N}(0,1)$，$U \sim \mathrm{Uniform}(0,1)$ 用以采样当日是否发生跳跃，$\xi_t \sim \mathcal{N}(\mu^{\tau}_J, (\sigma^{\tau}_J)^2)$ 为对数跳跃幅度。当日跳跃概率 $\lambda^{\tau}\Delta t$ 已远小于 1，故采用伯努利近似（一日最多一次跳跃）。

数值稳定性方面，对模拟价格作软钳制 $S^{\tau}_t \in [0.2\, P^{(0)},\, 5\, P^{(0)}]$，避免长程模拟中出现指数爆炸。

\subsection{涨跌幅与 K 线生成}

K 线 OHLC 由日级递推自然生成：
\begin{align}
O_i &= S^{\tau}_{i-1} \\
C_i &= S^{\tau}_{i} \\
H_i &= \max(O_i, C_i) \cdot (1 + U^{H}_i),\quad U^{H}_i \sim \bigl|\mathcal{N}(0, (0.3\sigma^{\tau}/\sqrt{252})^2)\bigr| + 0.003 \\
L_i &= \min(O_i, C_i) \cdot (1 - U^{L}_i),\quad U^{L}_i \sim \bigl|\mathcal{N}(0, (0.3\sigma^{\tau}/\sqrt{252})^2)\bigr| + 0.003
\end{align}

成交量与日内波动幅度联动，模拟真实市场"放量伴随大幅波动"的现象：
\begin{equation}
V_i \sim \mathrm{Uniform}(5\!\times\!10^4, 3\!\times\!10^5) \cdot \Bigl(1 + 5 \cdot \frac{|C_i - O_i|}{O_i}\Bigr)
\end{equation}

涨跌幅定义为：
\begin{equation}
\Delta\%_i = \frac{C_i - O_i}{O_i} \times 100\%
\end{equation}

\subsection{交易成本}

买入或卖出时，总成本/收入为：
\begin{equation}
\mathrm{Total} = P_{\mathrm{trade}} \times Q
\end{equation}
其中 $P_{\mathrm{trade}}$ 为成交价，$Q$ 为数量（份）。余额更新规则：
\begin{equation}
B_{\mathrm{new}} = B_{\mathrm{old}} - \mathrm{Total} \quad (\text{买入}), \qquad B_{\mathrm{new}} = B_{\mathrm{old}} + \mathrm{Total} \quad (\text{卖出})
\end{equation}

\subsection{对前端的可视化暴露}

为让用户直观感知市场逻辑，第 \ref{sec:macro} 节定义的宏观因子向量 $\mathbf{F}_t$ 以 5 秒一拍的频率通过 WebSocket 通道（\texttt{/api/exchange/ws/macro}）推送至前端"市场宏观面板"，并以箭头标注与上一拍的方向变化。当 WebSocket 不可用时，前端降级为 60 秒 HTTP 轮询 \texttt{/api/exchange/macro}，保障可用性。

%==============================================================================
\section{定投 (DCA) 原理}
%==============================================================================

\subsection{定投公式}

定时定投 (Dollar Cost Averaging) 的核心思想是固定间隔、固定金额买入，以平滑价格波动。设：
\begin{itemize}
\item $A$：每次定投数量（份）
\item $T$：定投间隔（分钟）
\item $N$：计划执行次数
\item $P_1, P_2, \ldots, P_N$：各次执行时的价格
\end{itemize}

则总投入（以国脉币计）为：
\begin{equation}
\mathrm{Cost}_{\mathrm{total}} = A \cdot \sum_{i=1}^{N} P_i
\end{equation}

获得的总份额：
\begin{equation}
Q_{\mathrm{total}} = N \cdot A
\end{equation}

平均成本：
\begin{equation}
\bar{P} = \frac{\mathrm{Cost}_{\mathrm{total}}}{Q_{\mathrm{total}}} = \frac{1}{N} \sum_{i=1}^{N} P_i
\end{equation}

\subsection{与一次性投入对比}

若一次性在 $t=1$ 时全仓买入，成本为 $N \cdot A \cdot P_1$。定投的优越性体现在当价格下行时，$\bar{P} < P_1$，从而降低平均持仓成本。

%==============================================================================
\section{持仓与盈亏计算}
%==============================================================================

\subsection{持仓聚合}

设用户 $u$ 的交易记录集合为 $\mathcal{R}_u$，对标的 $\tau$ 的净持仓为：
\begin{equation}
Q_u(\tau) = \sum_{r \in \mathcal{R}_u, r.\mathrm{symbol}=\tau} \bigl( \mathbb{1}_{r.\mathrm{action}=\mathtt{buy}} \cdot r.\mathrm{amount} - \mathbb{1}_{r.\mathrm{action}=\mathtt{sell}} \cdot r.\mathrm{amount} \bigr)
\end{equation}

\subsection{成本与市值}

持仓成本（仅考虑买入）：
\begin{equation}
C_u(\tau) = \sum_{r: \mathrm{buy}} r.\mathrm{amount} \cdot r.\mathrm{price}
\end{equation}

平均成本：
\begin{equation}
\bar{C}_u(\tau) = \frac{C_u(\tau)}{Q_u(\tau)}
\end{equation}

当前市值与浮动盈亏：
\begin{equation}
V_u(\tau) = Q_u(\tau) \cdot P_{\mathrm{current}}(\tau), \quad \mathrm{Profit}_u(\tau) = V_u(\tau) - C_u(\tau)
\end{equation}

盈亏比例：
\begin{equation}
\mathrm{Profit}\% = \frac{P_{\mathrm{current}} - \bar{C}}{\bar{C}} \times 100\%
\end{equation}

%==============================================================================
\section{AI 工具调用原理}
%==============================================================================

\subsection{Function Calling 流程}

AI 助手采用大语言模型 (LLM) 的 Function Calling 能力，将用户自然语言意图映射为结构化工具调用。

设用户输入为 $m$，工具集合为 $\mathcal{T} = \{t_1, t_2, \ldots, t_k\}$，则：
\begin{equation}
\mathrm{LLM}(m \mid \mathcal{T}) \to \begin{cases}
\text{直接回复} & \text{若无需工具} \\
(t_i, \mathrm{args}) & \text{若需调用 } t_i
\end{cases}
\end{equation}

\subsection{工具定义}

各工具 $t_i$ 由名称、描述、参数 schema 组成。例如余额查询：
\begin{equation}
t_{\mathrm{balance}} = (\mathtt{get\_user\_balance}, \emptyset)
\end{equation}

市场分析：
\begin{equation}
t_{\mathrm{market}} = (\mathtt{get\_market\_analysis}, \{\mathtt{symbol}: \mathrm{optional}\})
\end{equation}

\subsection{多轮调用}

若 LLM 返回工具调用，则执行工具得到结果 $r$，将 $(t_i, \mathrm{args}, r)$ 作为新消息反馈给 LLM，直至 LLM 不再请求工具：
\begin{equation}
m_0 = m \to (t_1, a_1) \to r_1 \to m_1 \to (t_2, a_2) \to r_2 \to \cdots \to \text{最终回复}
\end{equation}

%==============================================================================
\section{操作教程}
%==============================================================================

本节提供 MetaBank 各功能模块的详细操作步骤，便于用户快速上手。

%----------------------------------------------------------------------------
\subsection{环境准备}
%----------------------------------------------------------------------------

\begin{enumerate}
\item 安装 Python 3.8+ 与 Node.js 18+
\item 克隆项目：\texttt{git clone <repo> \&\& cd metabank}
\item 后端：\texttt{cd backend \&\& pip install -r requirements.txt \&\& python main.py}
\item 前端：\texttt{cd frontend \&\& npm install \&\& npm run dev}
\item 浏览器访问 \texttt{http://localhost:5173}
\end{enumerate}

\textbf{默认账户：} 管理员 \texttt{admin} / \texttt{admin123}；新用户注册自动获得 10,000 国脉币 (GMC)。

\textbf{数据初始化：} 可执行 \texttt{python backend/scripts/seed\_fake\_users.py} 与 \texttt{python backend/scripts/seed\_trades\_orders.py} 生成演示数据。

%----------------------------------------------------------------------------
\subsection{注册与登录}
%----------------------------------------------------------------------------

\textbf{步骤 1：} 打开首页，点击右上角「注册」。

\textbf{步骤 2：} 填写用户名（如 \texttt{demo}）和密码（如 \texttt{123456}），点击「注册」。

\textbf{步骤 3：} 注册成功后自动登录并跳转到应用内，系统分配专属钱包地址，格式为 $\mathtt{0xGM}$ 开头的 40 位十六进制串。

\textbf{步骤 4：} 若已注册，在首页点击「登录」，输入用户名和密码即可进入。

%----------------------------------------------------------------------------
\subsection{仪表盘}
%----------------------------------------------------------------------------

登录后默认进入仪表盘，主要元素包括：

\begin{itemize}
\item \textbf{余额卡片：} 显示当前国脉币 (GMC) 余额与钱包地址，可点击复制。
\item \textbf{市场概览：} 展示 GM、GMAI、GMC、GMFT、METAV、AIFIN 等标的的实时价格与涨跌幅 $\Delta\%$。
\item \textbf{快捷操作：} 转账、交易所、商城等入口。
\item \textbf{最近交易：} 最近资金变动记录。
\end{itemize}

%----------------------------------------------------------------------------
\subsection{钱包与转账}
%----------------------------------------------------------------------------

\textbf{查看资产：} 点击左侧「钱包」，可查看完整交易流水与余额。

\textbf{转账操作：}
\begin{enumerate}
\item 点击「转账」按钮。
\item 输入对方钱包地址（可从其他用户处获取，或使用 \texttt{admin} 用户地址测试）。
\item 输入转账金额（单位：GMC）。
\item 点击「确认」完成转账，交易即时到账。
\end{enumerate}

\textbf{注意：} 平台内所有金额均为国脉币 GMC，是平台虚拟代币，非人民币。

%----------------------------------------------------------------------------
\subsection{商城购物}
%----------------------------------------------------------------------------

\textbf{步骤 1：} 点击左侧「商城」，浏览商品列表。商品分类包括虚拟商品（AI 月卡、虚拟土地、NFT）、实物商品（智能手环、体检套餐）等。

\textbf{步骤 2：} 点击目标商品的「购买」按钮。

\textbf{步骤 3：} 选择购买数量，实物商品需填写收货地址与手机号。

\textbf{步骤 4：} 确认购买，系统扣除国脉币并生成订单，可在「我的订单」中查看。

%----------------------------------------------------------------------------
\subsection{交易所交易}
%----------------------------------------------------------------------------

\textbf{步骤 1：} 点击左侧「交易所」，上方为 K 线图，下方为交易表单。

\textbf{步骤 2：} 选择「买入」或「卖出」，选择标的（如国脉 AI 概念 GMAI）。

\textbf{步骤 3：} 输入数量（份），系统显示预估金额 $P \times Q$。

\textbf{步骤 4：} 点击「买入」或「卖出」，成交后余额与持仓自动更新。

\textbf{步骤 5：} 在「持仓」与「交易历史」中查看投资组合与盈亏。盈亏比例由公式 $\mathrm{Profit}\% = (P_{\mathrm{current}} - \bar{C})/\bar{C} \times 100\%$ 计算。

%----------------------------------------------------------------------------
\subsection{社区动态}
%----------------------------------------------------------------------------

\textbf{步骤 1：} 点击左侧「社区」，切换到「社区动态」标签。

\textbf{步骤 2：} 浏览其他用户发布的动态，包括生活分享、运动打卡、商城购物、交易理财等。

\textbf{步骤 3：} 可对动态进行「点赞」或「评论」。

\textbf{步骤 4：} 切换到「消息」标签可与好友私聊，支持在聊天中直接转账国脉币。

%----------------------------------------------------------------------------
\subsection{AI 智能助手}
%----------------------------------------------------------------------------

AI 助手支持自然语言交互与工具调用，常用操作如下。

\textbf{查看余额：} 输入「查看我的余额」或点击快捷按钮，AI 调用 \texttt{get\_user\_balance} 并展示结果。

\textbf{市场分析：} 输入「分析今日行情」或「GMAI 走势如何」，AI 调用 \texttt{get\_market\_analysis} 返回各标的行情与建议。

\textbf{查用户动态：} 输入「老张头最近在干嘛」，AI 调用 \texttt{get\_user\_activity} 汇总该用户最近动态。

\textbf{买入标的：} 输入「买入 50 份 GMAI」，AI 调用 \texttt{buy\_stock} 执行交易。

\textbf{设置定投：} 输入「设置 GMAI 定投，每次 10 份，每 60 分钟」，AI 调用 \texttt{set\_dca\_task} 创建定投任务。

\textbf{注意：} 所有金额均以国脉币 GMC 为单位。

%----------------------------------------------------------------------------
\subsection{定投任务管理}
%----------------------------------------------------------------------------

定投任务由 AI 助手或后台创建，系统按设定间隔自动执行买入。用户可在 AI 助手中通过自然语言查询或修改定投计划。

定投参数包括：标的符号、每次数量 $A$、间隔时间 $T$（分钟）、执行次数 $N$。平均成本满足 $\bar{P} = \frac{1}{N}\sum_{i=1}^{N} P_i$。

%----------------------------------------------------------------------------
\subsection{管理后台}
%----------------------------------------------------------------------------

\textbf{步骤 1：} 使用管理员账户 \texttt{admin} / \texttt{admin123} 登录。

\textbf{步骤 2：} 登录后左侧出现「管理后台」入口，点击进入。

\textbf{步骤 3：} 后台功能包括：
\begin{itemize}
\item \textbf{概览：} 用户数、订单数、交易量等统计。
\item \textbf{用户管理：} 查看用户列表，可调整用户余额。
\item \textbf{订单管理：} 处理商城订单状态。
\item \textbf{商品管理：} 添加、编辑、删除商品。
\item \textbf{系统设置：} 配置 AI 大模型 API Key 等。
\end{itemize}

%----------------------------------------------------------------------------
\subsection{API 调用}
%----------------------------------------------------------------------------

后端默认运行于 \texttt{http://localhost:8000}，除登录、注册等公开接口外，其余接口需在请求头中携带 JWT：
\begin{equation}
\mathtt{Authorization}: \mathtt{Bearer}\;\; \mathrm{JWT}
\end{equation}
登录接口 \texttt{POST /api/auth/login} 返回的 \texttt{access\_token} 即为 JWT。

%==============================================================================
\section{附录}
%==============================================================================

\subsection{符号表（LaTeX 数学符号）}

本文档中的公式与符号均使用 LaTeX 数学模式编写，下表列出常用符号含义：

\begin{table}[h]
\centering
\begin{tabular}{cl}
\toprule
符号 & 含义 \\
\midrule
$\mathcal{S}, \mathcal{F}, \mathcal{B}, \mathcal{D}, \mathcal{A}$ & 集合（系统、前端、后端、数据、AI） \\
$P^{(0)}$ & 基准价格 \\
$P_t$ & 时刻 $t$ 的价格 \\
$\varepsilon, \sigma$ & 希腊字母（扰动、标准差） \\
$\mathcal{N}(\mu, \sigma^2)$ & 正态分布 \\
$B$ & 用户余额 \\
$Q$ & 持仓数量 \\
$\bar{C}, \bar{P}$ & 平均成本、平均价格 \\
$\mathtt{GMC}$ & 国脉币（等宽字体表示代码/常量） \\
$\land, \implies$ & 逻辑与、蕴含 \\
$\|$ & 字符串连接 \\
\bottomrule
\end{tabular}
\end{table}

\subsection{文档技术说明}

本手册使用 \texttt{article} 文档类，依赖 \texttt{xeCJK}、\texttt{amsmath}、\texttt{hyperref} 等宏包，通过 XeLaTeX 编译以正确渲染中文与数学公式。


\end{document}
